
\section{Method}

\subsection{Overview}

NHA runs two attention streams in parallel and combines them via a token-level 
gate:

\begin{enumerate}
\item \textbf{Local stream} $O^\text{local}$: sliding-window attention.
\item \textbf{Global stream} $O^\text{global}$: RFF linear attention.
\item \textbf{Gate} $g$: per-token sigmoid.
\item \textbf{Fusion} $O = g \odot O^\text{local} + (1 - g) \odot O^\text{global}$.

\end{enumerate}
\textbf{Figure 4.1} illustrates the pipeline.

\begin{lstlisting}
        Q, K, V $\in$ R^{B $\times$ H $\times$ L $\times$ d\_h}
                 |
        +--------+--------+
                           
   +----------+      +----------+
   |  Local   |      |  Global  |
   |  window  |      |   RFF    |
   | O\_local  |      | O\_global |
   +----+-----+      +----+-----+
        |                 |
        +--------+--------+
                  
        +-----------------+
        |  Gate g(x)      |
        |  g O\_l+(1-g) O\_g|
        +--------+--------+
                  
              out\_proj
\end{lstlisting}

\subsection{Local Window Stream}

\subsubsection{Definition}

For window size $w$:

\[
S\_{ij} = \frac{Q\_i^\top K\_j}{\sqrt{d\_h}}, \quad
\text{mask}\_{ij} = \mathbb{1}[|i - j| \leq w]
\tag{4.1}
\]

\[
O\_i^\text{local} = \text{softmax}(S \odot \text{mask})\_i \cdot V
\tag{4.2}
\]

\subsubsection{Chunked Computation}

To bound memory, we process the query dimension in chunks of 64:

\begin{lstlisting}
for i in range(0, N, chunk):
    end = min(i + chunk, N)
    q\_block = Q[:, :, i:end]
    k\_start = max(0, i - window)
    k\_end = min(N, end + window)
    k\_block = K[:, :, k\_start:k\_end]
    # ... compute scores, mask, softmax ...
\end{lstlisting}

This gives $\mathcal{O}(B H C (2w + C))$ peak memory instead of 
$\mathcal{O}(B H L^2)$.

\subsubsection{Causal Masking}

For decoders:

\[
\text{mask}\_{ij}^\text{causal} = \mathbb{1}[j \leq i \text{ and } i - j \leq w]
\tag{4.3}
\]

\subsection{Global RFF Stream}

\subsubsection{RFF Projection}

We build $\Omega \in \mathbb{R}^{d_h \times m}$ via block QR:

\begin{lstlisting}
blocks = []
remaining = m
while remaining \textgreater{} 0:
    block\_size = min(d\_h, remaining)
    g = torch.randn(d\_h, block\_size)
    q, \_ = torch.linalg.qr(g)
    blocks.append(q)
    remaining -= block\_size
W = torch.cat(blocks, dim=1) * math.sqrt(d\_h)
\end{lstlisting}

\textbf{Orthogonality.} Each block $Q_i$ has orthonormal columns. The scale 
$\sqrt{d_h}$ matches the expected variance of Gaussian projections.

\subsubsection{Feature Map}

\[
\phi(x) = \frac{1}{\sqrt{m}} \left(\mathrm{elu}(\Omega^\top x) + 1\right)
\tag{4.4}
\]

\textbf{Properties:}
\begin{itemize}
\item \textbf{Non-negative:} $\mathrm{elu}(x) > -1$ $\to$ $\phi > 0$.
\item \textbf{Bounded:} $\phi \in [0, \infty)$.
\item \textbf{Smooth:} ELU is $C^1$.

\end{itemize}
\subsubsection{Non-Causal Computation}

For encoders:

\[
KV = \phi(K)^\top V \in \mathbb{R}^{m \times d\_h}, \quad
Z = \sum\_j \phi(K\_j) \in \mathbb{R}^m
\tag{4.5}
\]

\[
O\_i^\text{global} = \frac{\phi(Q\_i)^\top KV}{\phi(Q\_i)^\top Z + \epsilon}
\tag{4.6}
\]

\subsubsection{Causal Prefix Scan (Chunked)}

For decoders, we use a chunked prefix scan with cumulative state:

\[
S \in \mathbb{R}^{B \times H \times m \times d\_h}, \quad
Z \in \mathbb{R}^{B \times H \times m}
\]

Initialize $S = 0$, $Z = 0$. For each chunk $[start, end)$:

\begin{enumerate}
\item Compute intra-chunk causal scores:
  $$
  A^\text{intra} = \phi(Q\_c) \phi(K\_c)^\top \odot \text{tril}(1)
  \tag{4.7}
  $$

\item Compute inter-chunk contribution:
  $$
  O^\text{inter} = \phi(Q\_c) S
  \tag{4.8}
  $$

\item Combine:
  $$
  O\_c = \frac{A^\text{intra} V\_c + O^\text{inter}}{(A^\text{intra} \cdot 1 + \phi(Q\_c) Z) + \epsilon}
  \tag{4.9}
  $$

\item Update state:
  $$
  S \leftarrow S + \phi(K\_c)^\top V\_c, \quad Z \leftarrow Z + \sum \phi(K\_c)
  \tag{4.10}
  $$

\end{enumerate}
\textbf{Complexity.} $\mathcal{O}(B H L m d_h)$ time, $\mathcal{O}(B H m d_h)$ memory.

\subsubsection{Numerical Stability}

We compute the RFF stream in \textbf{FP32} inside an autocast-disabled block:

\begin{lstlisting}
with torch.autocast(device\_type=..., enabled=False):
    phi\_Q\_f = phi\_Q.float()
    phi\_K\_f = phi\_K.float()
    V\_f = V\_masked.float()
    # ... computation ...
\end{lstlisting}

This avoids FP16 overflow for $L > 512$.

\subsection{Token-Level Gate}

\subsubsection{Parameterization}

\[
g\_i = \sigma(w^\top x\_i + b)
\tag{4.11}
\]

where $x_i$ is the input hidden state at position $i$, and $w \in \mathbb{R}^d$, 
$b \in \mathbb{R}$.

\subsubsection{Initialization}

We initialize $w = 0$ and $b = 3.0$, so:

\[
g\_i = \sigma(0 + 3) = \sigma(3) \approx 0.95 \quad \text{for all } i
\tag{4.12}
\]

\textbf{Why bias toward local?} The local window matches the pretrained model's 
short-range behavior. Starting global-leaning would immediately degrade 
accuracy.

\subsubsection{Per-Token vs Per-Head}

NHA uses \textbf{per-token} gating (single scalar per position), while HALO-ReGLA 
uses per-head. Per-token is:

\begin{itemize}
\item \textbf{Parameter-efficient:} $d + 1$ parameters instead of $d H + H$.
\item \textbf{Position-aware:} Can adapt to local vs global needs per position.
\item \textbf{Simpler:} No head dimension to manage.

\end{itemize}
Empirically, per-token works as well as per-head for our tasks.

\subsection{Complete Forward Pass}

\textbf{Algorithm 1:} NHA forward
\begin{lstlisting}
Input:  hidden\_states $\in$ R^{B $\times$ L $\times$ d}
        attention\_mask, position\_embeddings, is\_causal
        window\_size w, rff\_features m
Output: output $\in$ R^{B $\times$ L $\times$ d}

 1: Q, K, V $\leftarrow$ proj(hidden\_states)
 2: Q, K $\leftarrow$ apply\_rope(Q, K, position\_embeddings)
 3: K, V $\leftarrow$ repeat\_kv(K, V, H / H\_kv) if GQA
 4:
 5: # Local stream (chunked)
 6: O\_local $\leftarrow$ sliding\_window\_attention(Q, K, V, w, is\_causal)
 7:
 8: # Global RFF stream
 9:  \_Q $\leftarrow$ elu($\Omega$^T Q) + 1
10:  \_K $\leftarrow$ elu($\Omega$^T K) + 1
11: if is\_causal:
12:     (S, Z) $\leftarrow$ initialize\_cumulative\_state()
13:     for chunk in chunks(L, 64):
14:         O\_global[chunk] $\leftarrow$ RFF\_chunk( \_Q,  \_K, V, S, Z)
15:         (S, Z) $\leftarrow$ update\_state(S, Z,  \_K[chunk], V[chunk])
16: else:
17:     KV $\leftarrow$  \_K^T V
18:     Z $\leftarrow$ sum( \_K)
19:     O\_global $\leftarrow$  \_Q @ KV / ( \_Q @ Z + $\epsilon$)
20:
21: # Token-level gate
22: g $\leftarrow$ sigmoid(gate(hidden\_states))         # (B, L, 1)
23: O $\leftarrow$ g * O\_local + (1 - g) * O\_global
24: O $\leftarrow$ O * valid\_mask
25: O $\leftarrow$ reshape and out\_proj
26: return O
\end{lstlisting}

\subsection{Complexity Analysis}

\begin{table}[H]
\centering
\small
\begin{tabular}{llll}
\toprule
\textbf{Component} & \textbf{Local} & \textbf{Global RFF} & \textbf{NHA (total)} \\
\midrule
Time & $\mathcal{O}(L w d_h)$ & $\mathcal{O}(L m d_h)$ & $\mathcal{O}(L (w + m) d_h)$ \\
Memory & $\mathcal{O}(L w)$ & $\mathcal{O}(m d_h)$ & $\mathcal{O}(L w + m d_h)$ \\
vs Dense & $\mathcal{O}(L^2 d_h)$ & -- & \textbf{$\mathcal{O}(L (w + m) d_h)$} \\
\bottomrule
\end{tabular}
\end{table}

\textbf{Key insight.} NHA is \textbf{linear} in $L$ for both branches. This is the 
primary efficiency benefit.

\subsection{Parameter Analysis}

\begin{table}[H]
\centering
\small
\begin{tabular}{lll}
\toprule
\textbf{Component} & \textbf{Baseline} & \textbf{NHA} \\
\midrule
$W_Q$ & $d H d_h$ & Same \\
$W_K$ & $d H_{kv} d_h$ & Same \\
$W_V$ & $d H_{kv} d_h$ & Same \\
$W_O$ & $H d_h d$ & Same \\
\textbf{Gate} & 0 & $d + 1$ \\
\textbf{RFF proj} & 0 & $d_h m$ (buffer, not trainable) \\
LoRA & 0 & $2 r (\cdots)$ \\
\bottomrule
\end{tabular}
\end{table}

\textbf{Total: same as MHA + $d + 1$ parameters (gate) + RFF buffer.}

\subsection{Theoretical Notes}

\subsubsection{Kernel Approximation}

With orthogonal $\Omega$ and $m \to \infty$:

\[
\phi(q)^\top \phi(k) \xrightarrow{m \to \infty} \mathbb{E}\_\omega[\ldots] = k\_\text{exp}(q, k)
\tag{4.13}
\]

where $k_\text{exp}$ is the exp kernel.

\subsubsection{Gate Convergence}

The gate converges to a stationary point where:

\[
\frac{\partial \mathcal{L}}{\partial g\_i} = (\nabla\_O \mathcal{L}) \cdot (O^\text{local}\_i - O^\text{global}\_i) = 0
\tag{4.14}
\]

or hits the boundary ($g_i \to 0$ or $g_i \to 1$).

\subsubsection{Causality}

For causal models, the chunked prefix scan ensures:

\[
O\_i^\text{global} = f(\phi(Q\_j), \phi(K\_j), V\_j : j \leq i)
\tag{4.15}
\]

No future information leaks.
\begin{lstlisting}

---
\end{lstlisting}